\documentclass[12pt]{article}
\usepackage{notes-project}

\begin{document}

\begin{center}
  {\LARGE\bfseries Lecture 4: Limit of Sequence}\\[0.25em]
  {\large Mathematical Analysis}\\[0.15em]
\end{center}

\section{Convergent Sequences} 

\begin{dfn}[Convergent Sequence]{convergent-sequence}
  A sequence $a: \NNOne \to F$, where $F$ is a \xref[ordered field]{lec01::def:ordered-field}, is said to be \emph{convergent} if
  \[
    \exists \ell \in F,\, \forall \varepsilon \in \apos{F},\, \exists N \in \NNOne,\, \forall n \in \NNOne,\, n \geq N,\, \dist{a_n}{\ell} < \varepsilon
  \]
\end{dfn}

\begin{rem}[The $\varepsilon$-$N$ Challenge]{epsilon-n-challenge}
A value $\ell$ passes the $\varepsilon$-$N$ challenge if for any proposed $\varepsilon > 0$, there exists an index $N \in \NNOne$ such that for all $n \in \NNOne$ with $n \geq N$, $\dist{a_n}{\ell} < \varepsilon$
\end{rem}

\begin{rem}
In the \xref[The $\varepsilon$-$N$ Challenge]{rem:epsilon-n-challenge}, the value of $N$ may depend on the choice of $\varepsilon$. And the value $N$ may be not unique (ie, $N' > N$ also works)
\end{rem}

\begin{thm}[Uniqueness of Limit]{limit-uniqueness}
  Let $F$ be an \xref[ordered field]{lec01::def:ordered-field}, $a : \NNOne \to F$, and $\ell_1, \ell_2 \in F$. Then
  \[
    \bigl(\forall \varepsilon \in \apos{F},\, \exists N \in \NNOne,\, \forall n \in \NNOne,\, n \geq N,\, \dist{a_n}{\ell_1} < \varepsilon\bigr)
  \]\[
    \land
  \]\[
    \bigl(\forall \varepsilon \in \apos{F},\, \exists N \in \NNOne,\, \forall n \in \NNOne,\, n \geq N,\, \dist{a_n}{\ell_2} < \varepsilon\bigr)
  \]\[
    \implies \ell_1 = \ell_2
  \]
\end{thm}

\begin{proof}
Assume for contradiction that $\ell_1 \neq \ell_2$. Set
\[
  \varepsilon = \frac{\dist{\ell_1}{\ell_2}}{2}.
\]
Since $\ell_1 \neq \ell_2$, we have $\varepsilon > 0$

By the hypothesis for $\ell_1$, there exists $N_1 \in \NNOne$ such that for all $n \geq N_1$, $\dist{a_n}{\ell_1} < \varepsilon$. Similarly, there exists $N_2 \in \NNOne$ such that for all $n \geq N_2$, $\dist{a_n}{\ell_2} < \varepsilon$

Let $N = \max(N_1, N_2)$. For any $n \geq N$, both hypotheses hold:
\[
  \dist{a_n}{\ell_1} < \varepsilon \quad\text{and}\quad \dist{a_n}{\ell_2} < \varepsilon
\]

By the \xref[triangle inequality for distance]{lec03::lem:triangle-inequality-distance} applied via the intermediate point $a_n$,
\[
  \dist{\ell_1}{\ell_2} \leq \dist{\ell_1}{a_n} + \dist{a_n}{\ell_2} = \dist{a_n}{\ell_1} + \dist{a_n}{\ell_2},
\]
where the last equality uses \xref[symmetry of distance]{lec03::lem:distance-symmetry}. Combining,
\[
  \dist{\ell_1}{\ell_2} \leq \dist{a_n}{\ell_1} + \dist{a_n}{\ell_2} < 2\varepsilon = \dist{\ell_1}{\ell_2},
\]
a contradiction. Hence $\ell_1 = \ell_2$
\end{proof}

\begin{rem}[Compact Form Using $\exists!$]{limit-uniqueness-compact}
\missingprereq{Combined with the existence from the \xref[definition of convergent sequence]{def:convergent-sequence}, this theorem states: if $a$ is convergent, then there exists a unique such $\ell$:
\[
  \exists! \ell \in F,\, \forall \varepsilon \in \apos{F},\, \exists N \in \NNOne,\, \forall n \in \NNOne,\, n \geq N,\, \dist{a_n}{\ell} < \varepsilon.
\]
Here $\exists!$ is shorthand for ``existence and uniqueness''; it should be formally introduced in \texttt{lec01} as a FOL convention before being used}
\end{rem}

\section{Convergence Challenge Functions}

\begin{dfn}[Convergence Challenge Function]{convergence-challenge-function}
  A function $g : \RR \to \RR$ is a \emph{convergence challenge function} if it satisfies two conditions:
  \begin{itemize}[nosep]
    \item \emph{Positivity on positives}: $\forall \gamma \in \apos{\RR},\, g(\gamma) > 0$
    \item \emph{Arbitrarily small values}: $\forall \delta \in \apos{\RR},\, \exists \gamma \in \apos{\RR},\, g(\gamma) < \delta$
  \end{itemize}
\end{dfn}

\begin{rem}[Intuition: Covering Arbitrarily Small Positives]{challenge-intuition}
\end{rem}

\begin{thm}[Convergence Challenge Replacement]{convergence-challenge-replacement}
  Let \(g\) be a
  \xref[convergence challenge function]{def:convergence-challenge-function},
  and let \(a:\NNOne\to X\) be a sequence in a
  \xref[metric space]{lec03::def:metric-space} \((X,d)\). Then
  \[
    \bigl(\exists \ell \in X,\, \forall \varepsilon \in \apos{\RR},\, \exists N \in \NNOne,\, \forall n \in \NNOne,\, n \geq N,\, d(\term{a}{n},\ell) < \varepsilon\bigr)
  \]\[
    \iff
  \]\[
    \bigl(\exists \ell \in X,\, \forall \varepsilon \in \apos{\RR},\, \exists N \in \NNOne,\, \forall n \in \NNOne,\, n \geq N,\, d(\term{a}{n},\ell) < g(\varepsilon)\bigr)
  \]
\end{thm}

\begin{proof}
  Fix the witness $\ell \in X$ in either direction

  \noindent\textbf{($\Rightarrow$)} For $\varepsilon > 0$, the positivity-on-positives condition gives $g(\varepsilon) > 0$. Apply the left-hand condition at the tolerance $g(\varepsilon)$

  \noindent\textbf{($\Leftarrow$)} For $\varepsilon > 0$, the \xref[challenge-function property of $g$]{def:convergence-challenge-function} gives $\gamma > 0$ with $g(\gamma) < \varepsilon$. The hypothesis at $\gamma$ gives $N$ such that
  \[
    d(\term{a}{n},\ell) < g(\gamma) < \varepsilon \quad \text{for all } n \geq N
  \]
\end{proof}

\begin{cor}[Equality via Challenge Functions]{equality-challenge-replacement}
  Let \(g\) be a
  \xref[convergence challenge function]{def:convergence-challenge-function},
  and let \(x,y\in X\) in a
  \xref[metric space]{lec03::def:metric-space} \((X,d)\). Then
  \[
    x=y\iff\forall\varepsilon\in\apos{\RR},\quad d(x,y)<g(\varepsilon)
  \]
\end{cor}

\begin{proof}
  \noindent\textbf{(\(\Rightarrow\))}
  If \(x=y\), then \(d(x,y)=0<g(\varepsilon)\) for every
  \(\varepsilon>0\)

  \noindent\textbf{(\(\Leftarrow\))}
  If \(x\neq y\), then \(d(x,y)>0\). The challenge-function property
  gives \(\gamma>0\) with \(g(\gamma)<d(x,y)\), contradicting
  \(d(x,y)<g(\gamma)\). Hence \(x=y\)
\end{proof}

\begin{lem}[Restricting to Small Tolerances]{restricting-to-small-tolerances}
  Fix $r\in\apos{\RR}$. The function
  \[
    g:\RR\to\RR\defeq\gamma\map\min(\gamma,r/2)
  \]
  is a \xref[convergence challenge function]{def:convergence-challenge-function}
\end{lem}

\begin{proof}
  For every $\gamma>0$, $0<g(\gamma)\leq r/2<r$.
  For every $\delta>0$, choose $\gamma\defeq\delta/2>0$; then
  $g(\gamma)\leq\delta/2<\delta$
\end{proof}

\begin{rem}[Restricting the Tolerance]{restricting-the-tolerance}
  Since $0<g(\gamma)<r$ for $\gamma>0$,
  \xref[challenge replacement]{thm:convergence-challenge-replacement} lets us
  check $\forall\,0<\varepsilon<r$ instead of $\forall\varepsilon>0$,
  for any fixed $r>0$, eg, $\forall\,0<\varepsilon<1$
\end{rem}

\begin{lem}[Linear Scaling]{linear-is-challenge}
  For any $c \in \RR$ with $c > 0$, the function $g:\RR\to\RR\defeq\gamma\map c\cdot\gamma$ is a \xref[convergence challenge function]{def:convergence-challenge-function}
\end{lem}

\begin{proof}
  We verify both conditions

  \noindent\emph{Positivity on positives.} For any $\gamma \in \RR$ with $\gamma > 0$, both $c > 0$ and $\gamma > 0$, so $g(\gamma) = c\cdot\gamma > 0$

  \noindent\emph{Arbitrarily small values.} Fix any $\delta \in \RR$ with $\delta > 0$. We seek $\gamma \in \RR$ with $\gamma > 0$ and $g(\gamma) < \delta$

  Take $\gamma \defeq \delta / (c + 1)$, which is derived in \xref{rem:linear-scaling-witness-derivation}. Since $\delta > 0$ and $c + 1 > 0$ (because $c > 0$), we have $\gamma > 0$

  Now $g(\gamma) = c\cdot\gamma = c\cdot\delta / (c + 1)$. Since $c < c + 1$, we have $c / (c + 1) < 1$, hence $c\cdot\delta / (c + 1) < \delta$
\end{proof}

\begin{rem}[Deriving the Witness for Linear Scaling]{linear-scaling-witness-derivation}
  To find the witness $\gamma$ used in the proof of \xref{lem:linear-is-challenge}, we start from the target inequality $g(\gamma) < \delta$ and manipulate it to isolate $\gamma$:
  \[
    g(\gamma) < \delta
    \iff c\cdot\gamma < \delta
    \iff \gamma < \frac{\delta}{c}.
  \]
  Any $\gamma$ in the open interval $\left(0, \tfrac{\delta}{c}\right)$ will therefore satisfy $g(\gamma) < \delta$. To pick a concrete such $\gamma$, we use the fact that $c + 1 > c$, which gives
  \[
    \frac{\delta}{c + 1} < \frac{\delta}{c}.
  \]
  Hence we may take $\gamma \defeq \delta / (c + 1)$ as the witness
\end{rem}

\begin{rem}[Why $c \leq 0$ is Excluded]{linear-scaling-cnonpos}
  If $c \leq 0$, then for any $\gamma \in \RR$ with $\gamma > 0$, $g(\gamma) = c\cdot\gamma \leq 0$. This violates the positivity-on-positives condition, so $g$ is not a \xref[convergence challenge function]{def:convergence-challenge-function}
\end{rem}

\begin{lem}[Power]{power-is-challenge}
  For any positive integer $k$, the function $g:\RR\to\RR\defeq\gamma\map\gamma^k$ is a \xref[convergence challenge function]{def:convergence-challenge-function}
\end{lem}

\begin{proof}
  Both conditions follow by \xref[induction on $k$]{lec01::thm:natural-number-induction}. Positivity: $\gamma^k > 0$ for $\gamma > 0$. For arbitrarily small values, take $\gamma \defeq \min(\delta / 2,\, 1 / 2)$, so $\gamma > 0$ and $\gamma \leq 1$; induction gives $\gamma^k \leq \gamma \leq \delta / 2 < \delta$
\end{proof}

\begin{lem}[Summation]{sum-is-challenge}
  Let $f, g$ be \xref[convergence challenge functions]{def:convergence-challenge-function}, and suppose both are monotone increasing (i.e.\ $x \leq y \implies f(x) \leq f(y)$, and similarly for $g$). Then $f + g$, defined by $(f + g)(\gamma) \defeq f(\gamma) + g(\gamma)$, is a \xref[convergence challenge function]{def:convergence-challenge-function}
\end{lem}

\begin{proof}
  Positivity is immediate from $f(\gamma), g(\gamma) > 0$. For $\delta > 0$, pick $\gamma_1, \gamma_2 > 0$ with $f(\gamma_1), g(\gamma_2) < \delta / 2$; set $\gamma \defeq \min(\gamma_1, \gamma_2)$. Monotonicity gives $(f + g)(\gamma) \leq f(\gamma_1) + g(\gamma_2) < \delta$
\end{proof}

\begin{lem}[Composition]{challenge-composition}
  Let $f, g$ be \xref[convergence challenge functions]{def:convergence-challenge-function}, and suppose $f$ is monotone increasing (i.e.\ $x \leq y \implies f(x) \leq f(y)$). Then $f \circ g$ is a \xref[convergence challenge function]{def:convergence-challenge-function}
\end{lem}

\begin{proof}
  Positivity: $f(g(\gamma)) > 0$ since $g(\gamma) > 0$ and $f$ preserves positivity. For $\delta > 0$, pick $\eta > 0$ with $f(\eta) < \delta$, then $\gamma > 0$ with $g(\gamma) < \eta$. Monotonicity of $f$ gives $(f \circ g)(\gamma) \leq f(\eta) < \delta$
\end{proof}

\begin{thm}[Closed Form of Convergence]{convergence-closed-form}
  Let $F$ be an \xref[ordered field]{lec01::def:ordered-field}, let $a : \NNOne \to F$. Then
  \[
    \bigl(\exists \ell \in F,\, \forall \varepsilon \in \apos{F},\, \exists N \in \NNOne,\, \forall n \in \NNOne,\, n \geq N,\, \dist{a_n}{\ell} < \varepsilon\bigr)
  \]\[
    \iff
  \]\[
    \bigl(\exists \ell \in F,\, \forall \varepsilon \in \apos{F},\, \exists N \in \NNOne,\, \forall n \in \NNOne,\, n \geq N,\, \dist{a_n}{\ell} \leq \varepsilon\bigr)
  \]
\end{thm}

\begin{proof}
  \noindent\textbf{($\Rightarrow$)} If $\dist{a_n}{\ell} < \varepsilon$, then $\dist{a_n}{\ell} \leq \varepsilon$

  \noindent\textbf{($\Leftarrow$)} Fix the witness \(\ell\in F\).
  For any \(\varepsilon\in\apos{F}\), apply the right-hand condition at
  \(\varepsilon/2>0\) to obtain \(N\in\NNOne\) such that
  \[
    n\geq N\implies\dist{a_n}{\ell}\leq\varepsilon/2<\varepsilon
  \]
\end{proof}

\section{Limit of a Sequence}

\begin{dfn}[Limit of a Sequence]{limit}
  Let $F$ be an \xref[ordered field]{lec01::def:ordered-field}. We say the sequence $a: \NNOne \to F$ \emph{has a limit of} $\ell \in F$ if
  \[
    \forall \varepsilon \in \apos{F},\, \exists N \in \NNOne,\, \forall n \in \NNOne,\, n \geq N,\, \dist{a_n}{\ell} < \varepsilon
  \]
\end{dfn}

\redirect[Converges To]{def}{converges-to}{def:limit}

\begin{dfn}[Set of Convergent Sequences]{convergent-sequences-set}
  Let $F$ be an \xref[ordered field]{lec01::def:ordered-field}. The \emph{set of convergent sequences in $F$} is
  \[
    \Conv(F) \defeq \{a : \NNOne \to F \mid a \text{ is \xref[convergent]{def:convergent-sequence}}\}
  \]
  \par\smallskip\noindent{\footnotesize\texttt{LaTeX: \detokenize{\Conv(F)}}}
\end{dfn}

\begin{dfn}[Limit Operator]{limit-op}
  Let $F$ be an \xref[ordered field]{lec01::def:ordered-field}. The \emph{limit operator} is the function
  \[
    \seqlim : \Conv(F) \to F \defeq a \map \text{the unique } \ell \in F \text{ such that } a \xref[\text{has a limit of}]{def:limit} \ell
  \]
  We write
  \[
    \seqlim(a) = \ell
  \]
  The input is the whole sequence, rather than an individual term
  \par\smallskip\noindent{\footnotesize\texttt{LaTeX: \detokenize{\seqlim(a)}}}
\end{dfn}

\redirect[Limit of a Sequence]{def}{limit-of-sequence}{def:limit-op}

\begin{rem}[Well-Definedness of the Limit]{limit-well-defined}
  The function $\seqlim$ in \xref{def:limit-op} is well-defined for two reasons:
  \begin{itemize}[nosep]
    \item \emph{Every input has an output}: for every $a \in \Conv(F)$, the \xref[definition of convergent sequence]{def:convergent-sequence} guarantees some $\ell \in F$ with $a \xref[\text{has a limit of}]{def:limit} \ell$
    \item \emph{Every input has exactly one output}: by \xref{thm:limit-uniqueness}, this $\ell$ is unique
  \end{itemize}
\end{rem} 

\begin{thm}[Addition of the Limit]{limit-addition}
  Let $F$ be an \xref[ordered field]{lec01::def:ordered-field}, and let $a, b \in \Conv(F)$. Then
  \[
    \seqlim(n\map\term{a}{n}+\term{b}{n}) = \seqlim(a)+\seqlim(b)
  \] 
\end{thm}

\begin{proof}
We want to prove
\[
  \begin{gathered}
    \forall\varepsilon\in \apos{F},\,\exists N\in\NNOne,\,
    \forall n\in\NNOne,\,n\geq N\implies\\
    \abs{\term{a}{n}+\term{b}{n}-\seqlim(a)-\seqlim(b)}<\varepsilon
  \end{gathered}
\]

According to the \xref[definition of a limit]{def:limit}, for every
\(\varepsilon\in \apos{F}\), there exist \(N_a,N_b\in\NNOne\) such that,
setting \(N_M\defeq\max(N_a,N_b)\), both estimates hold for \(n\geq N_M\):
\[
  \abs{\term{a}{n}-\seqlim(a)}<\varepsilon
  \quad\land\quad
  \abs{\term{b}{n}-\seqlim(b)}<\varepsilon
\]
Therefore, by the \xref[triangle inequality for distance]{lec03::lem:triangle-inequality-distance},
\[
  \begin{aligned}
    \abs{\term{a}{n}+\term{b}{n}-\seqlim(a)-\seqlim(b)}
    &\leq\abs{\term{a}{n}-\seqlim(a)}
      +\abs{\term{b}{n}-\seqlim(b)}\\
    &<\varepsilon+\varepsilon=2\cdot\varepsilon
  \end{aligned}
\]

For any \(\delta\in\apos{F}\), apply the preceding estimate with
\(\varepsilon=\delta/2>0\). The error is then less than
\(2\cdot(\delta/2)=\delta\), so
\[
  \seqlim(n\map\term{a}{n}+\term{b}{n})=\seqlim(a)+\seqlim(b)
\]
\end{proof}

\begin{thm}[Multiplication of the Limit]{limit-multiplication}
  Let $F$ be an \xref[ordered field]{lec01::def:ordered-field}, and let $a, b \in \Conv(F)$. Then
  \[
    \seqlim(n\map\term{a}{n}\cdot\term{b}{n})=\seqlim(a)\cdot\seqlim(b)
  \]
\end{thm}

\begin{proof}
For \(n\map\term{a}{n}\cdot\term{b}{n}\), we want to show that
\(\seqlim(a)\cdot\seqlim(b)\) is its limit, ie,
\[
  \begin{gathered}
    \forall\varepsilon\in \apos{F},\,\exists N\in\NNOne,\,
    \forall n\in\NNOne,\,n\geq N\implies\\
    \abs{\term{a}{n}\cdot\term{b}{n}-\seqlim(a)\cdot\seqlim(b)}<\varepsilon
  \end{gathered}
\]

According to the \xref[definition of a limit]{def:limit}, for every
\(\varepsilon\in \apos{F}\), there exist \(N_a,N_b\in\NNOne\) such that
\[
  \begin{aligned}
    n\geq N_a&\implies\abs{\term{a}{n}-\seqlim(a)}<\varepsilon,\\
    n\geq N_b&\implies\abs{\term{b}{n}-\seqlim(b)}<\varepsilon
  \end{aligned}
\]
Let \(N_M\defeq\max(N_a,N_b)\). Then, for every \(n\geq N_M\), both
inequalities hold

Adding and subtracting \(\seqlim(a)\cdot\term{b}{n}\), we have
\[
  \begin{aligned}
    &\abs{\term{a}{n}\cdot\term{b}{n}-\seqlim(a)\cdot\seqlim(b)}\\
    &\quad=\abs{\term{a}{n}\cdot\term{b}{n}
      -\seqlim(a)\cdot\term{b}{n}
      +\seqlim(a)\cdot\term{b}{n}-\seqlim(a)\cdot\seqlim(b)}\\
    &\quad\leq\abs{\term{a}{n}\cdot\term{b}{n}-\seqlim(a)\cdot\term{b}{n}}
      +\abs{\seqlim(a)\cdot\term{b}{n}-\seqlim(a)\cdot\seqlim(b)}\\
    &\quad=\abs{\term{b}{n}}\cdot\abs{\term{a}{n}-\seqlim(a)}
      +\abs{\seqlim(a)}\cdot\abs{\term{b}{n}-\seqlim(b)}
  \end{aligned}
\]
The inequality uses the
\xref[triangle inequality for distance]{lec03::lem:triangle-inequality-distance}

Therefore, when \(n\geq N_M\),
\[
  \begin{aligned}
    \abs{\term{a}{n}\cdot\term{b}{n}-\seqlim(a)\cdot\seqlim(b)}
    &\leq\abs{\term{b}{n}}\cdot\varepsilon+\abs{\seqlim(a)}\cdot\varepsilon\\
    &=\varepsilon\cdot\bigl(\abs{\term{b}{n}}+\abs{\seqlim(a)}\bigr)
  \end{aligned}
\]

Since all \xref[convergent sequences are bounded]{thm:convergent-sequences-are-bounded},
choose a strict bound \(M'\in F\) such that
\(\forall n\in\NNOne,\,\abs{\term{b}{n}}<M'\), and let
\(M\defeq\max(M',1)\). This bound is independent of \(\varepsilon\).
Hence \(M\geq1\) and
\[
  \abs{\term{a}{n}\cdot\term{b}{n}-\seqlim(a)\cdot\seqlim(b)}
  <\varepsilon\cdot\bigl(M+\abs{\seqlim(a)}\bigr)
\]

For any \(\delta\in\apos{F}\), apply the preceding estimate with
\[
  \varepsilon=\frac{\delta}{M+\abs{\seqlim(a)}}>0
\]
Since \(M\) is independent of the tolerance, the error is then less than
\(\delta\). Therefore,
\[
  \seqlim(n\map\term{a}{n}\cdot\term{b}{n})=\seqlim(a)\cdot\seqlim(b)
\]
\end{proof}

\begin{thm}[Scalar Multiplication of the Limit]{limit-scalar-multiplication}
  Let $F$ be an \xref[ordered field]{lec01::def:ordered-field}, let $a \in \Conv(F)$, and let $c \in F$. Then
  \[
    \seqlim(n\map c\cdot\term{a}{n})=c\cdot\seqlim(a)
  \]
\end{thm}

\begin{proof}
Define the constant sequence
\[
  p:\NNOne\to F\defeq n\map c
\]
For every \(\varepsilon\in \apos{F}\) and every \(n\in\NNOne\),
\[
  \abs{\term{p}{n}-c}=0<\varepsilon
\]
Hence \(p\in\Conv(F)\)
and \(\seqlim(p)=c\)

By the definition of \(p\),
\[
  \forall n\in\NNOne,\quad c\cdot\term{a}{n}=\term{p}{n}\cdot\term{a}{n}
\]
By \xref[multiplication of limits]{thm:limit-multiplication}, the sequence
\(n\map\term{p}{n}\cdot\term{a}{n}\) is convergent, and
\[
  \begin{aligned}
    \seqlim(n\map c\cdot\term{a}{n})
    &=\seqlim(n\map\term{p}{n}\cdot\term{a}{n})\\
    &=\seqlim(p)\cdot\seqlim(a)\\
    &=c\cdot\seqlim(a)
  \end{aligned}
\]
\end{proof}

\begin{thm}[Linearity of the Limit]{limit-linearity}
  Let $F$ be an \xref[ordered field]{lec01::def:ordered-field}, let $a, b \in \Conv(F)$, and let $c,d\in F$. Then
  \[
    \seqlim(n\map c\cdot\term{a}{n}+d\cdot\term{b}{n})
      =c\cdot\seqlim(a)+d\cdot\seqlim(b)
  \]
\end{thm}

\begin{proof}
This follows from \xref[addition of limits]{thm:limit-addition} and
\xref[scalar multiplication of limits]{thm:limit-scalar-multiplication}
\end{proof}

\begin{lem}[Reciprocal of the Limit]{limit-reciprocal}
  Let $F$ be an \xref[ordered field]{lec01::def:ordered-field}, and let
  $b\in\Conv(F)$ with $\forall n\in\NNOne,\,\term{b}{n}\neq0$
  and $\seqlim(b)\neq0$. Then
  \[
    \seqlim\left(n\map\frac{1}{\term{b}{n}}\right)
      =\frac{1}{\seqlim(b)}
  \]
\end{lem}

\begin{proof}
Define
\[
  p:\NNOne\to F\defeq n\map\frac{1}{\term{b}{n}}
\]
We check whether $1/\seqlim(b)$ is a limit of $p$

For $n\in\NNOne$,
\[
  \begin{aligned}
    \abs{\term{p}{n}-\frac{1}{\seqlim(b)}}
      &=\abs{\frac{1}{\term{b}{n}}-\frac{1}{\seqlim(b)}}\\
      &=\abs{\frac{\seqlim(b)-\term{b}{n}}
        {\term{b}{n}\cdot\seqlim(b)}}\\
      &=\frac{\abs{\seqlim(b)-\term{b}{n}}}
        {\abs{\term{b}{n}\cdot\seqlim(b)}}
  \end{aligned}
\]
To bound the denominator below, set $M\defeq\abs{\seqlim(b)}/2>0$.
By the \xref[convergence of $b$]{def:convergent-sequence}, choose
$N_0\in\NNOne$ such that
$n\geq N_0\implies\abs{\term{b}{n}-\seqlim(b)}<M$.
For $n\geq N_0$, the
\xref[triangle inequality]{lec03::lem:triangle-inequality-distance} gives
\[
  \abs{\seqlim(b)}
  \leq\abs{\seqlim(b)-\term{b}{n}}+\abs{\term{b}{n}}
  <M+\abs{\term{b}{n}}
\]
Since $\abs{\seqlim(b)}=2\cdot M$, we have $\abs{\term{b}{n}}>M$.
So
\[
  \exists M\in\apos{F},\,\exists N_0\in\NNOne,\,
  \forall n\in\NNOne,\,n\geq N_0\implies\abs{\term{b}{n}}>M
\]

Fix $\varepsilon\in\apos{F}$ and choose $N_b\in\NNOne$ so that
$n\geq N_b\implies\abs{\term{b}{n}-\seqlim(b)}<\varepsilon$.
Set $N\defeq\max(N_0,N_b)$. For $n\geq N$,
\[
  \abs{\term{b}{n}\cdot\seqlim(b)}>M\cdot\abs{\seqlim(b)},
  \qquad
  \abs{\seqlim(b)-\term{b}{n}}<\varepsilon
\]
Therefore,
\[
  \frac{\abs{\seqlim(b)-\term{b}{n}}}
    {\abs{\term{b}{n}\cdot\seqlim(b)}}
  <\frac{\varepsilon}{M\cdot\abs{\seqlim(b)}}
\]
For any \(\delta\in\apos{F}\), apply this estimate with
\[
  \varepsilon=\delta\cdot M\cdot\abs{\seqlim(b)}>0
\]
The error is then less than \(\delta\), so \(p\in\Conv(F)\) and
\[
  \seqlim(p)=\frac{1}{\seqlim(b)}
\]
\end{proof}

\begin{thm}[Division of the Limit]{limit-division}
  Let $F$ be an \xref[ordered field]{lec01::def:ordered-field}, and let $a, b \in \Conv(F)$ with $\forall n \in \NNOne, \term{b}{n} \neq 0$ and $\seqlim(b) \neq 0$. Then
  \[
    \seqlim\left(n\map\frac{\term{a}{n}}{\term{b}{n}}\right) = \frac{\seqlim(a)}{\seqlim(b)}
  \] 
\end{thm}

\begin{proof}
Define
\[
  p:\NNOne\to F\defeq n\map\frac{1}{\term{b}{n}}
\]
By the \xref[reciprocal lemma]{lem:limit-reciprocal},
$p\in\Conv(F)$ and $\seqlim(p)=1/\seqlim(b)$.
Since $\term{a}{n}/\term{b}{n}=\term{a}{n}\cdot\term{p}{n}$ for every $n\in\NNOne$,
\xref[multiplication of limits]{thm:limit-multiplication} gives
\[
  \begin{aligned}
    \seqlim\left(n\map\frac{\term{a}{n}}{\term{b}{n}}\right)
      &=\seqlim(n\map\term{a}{n}\cdot\term{p}{n})\\
      &=\seqlim(a)\cdot\seqlim(p)\\
      &=\frac{\seqlim(a)}{\seqlim(b)}
  \end{aligned}
\]
\end{proof}

\begin{thm}[Order Preservation of Limits]{limit-order-preservation}
  Let $F$ be an \xref[ordered field]{lec01::def:ordered-field}, and let $a, b \in \Conv(F)$ with $\forall n \in \NNOne,\, \term{a}{n} \leq \term{b}{n}$. Then
  \[
    \seqlim(a) \leq \seqlim(b)
  \]
\end{thm}

\begin{thm}[Sequence Characterization of Limit Points]{limit-point-sequence}
  Let \((X,d)\) be a \xref[metric space]{lec03::def:metric-space},
  let \(A\subseteq X\), and let \(x\in X\).
  The point \(x\) is a \xref[limit point]{lec03::def:limit-points} of \(A\)
  if and only if there exists a sequence \(a:\NNOne\to A\setminus\{x\}\)
  converging to \(x\), meaning
  \[
    \forall\varepsilon\in\apos{\RR},\,\exists N\in\NNOne,\,
    \forall n\in\NNOne,\quad
    n\geq N\implies d(\term{a}{n},x)<\varepsilon
  \]
\end{thm}

\begin{proof}
  According to definition of limit, \(x\) is a limit 
\end{proof}

\section{Bounded Sequences}

\begin{dfn}[Bounded Sequence]{bounded-sequence}
  Let \(F\) be an \xref[ordered field]{lec01::def:ordered-field}.
  A sequence \(a:\NNOne\to F\) is \emph{bounded} if
  \[
    \exists M\in F,\,\forall n\in\NNOne,\quad\dist{a(n)}{0}\leq M
  \]
\end{dfn} 

\begin{dfn}[Set of Bounded Sequences]{bounded-sequences-set}
  Let \(F\) be an \xref[ordered field]{lec01::def:ordered-field}.
  The set of bounded sequences in \(F\) is
  \[
    \Bdd(F)\defeq
    \{a:\NNOne\to F\mid
      a\text{ is \xref[bounded]{def:bounded-sequence}}\}
  \]
  \par\smallskip\noindent{\footnotesize\texttt{LaTeX: \detokenize{\Bdd(F)}}}
\end{dfn}

\begin{thm}[Convergent Sequences are Bounded]
  Let \(F\) be an \xref[ordered field]{lec01::def:ordered-field}
  and let \(a:\Conv(F)\). Then \(a\) is
  \xref[bounded]{def:bounded-sequence}, so \(a\in\Bdd(F)\)
\end{thm}

\begin{proof}
Suppose the convergent sequence $a$, $\seqlim a = \ell$

Take $\varepsilon = 1$, then there exists $N \in \NNOne$ such that for all $n \in \NNOne$ with $n > N$, $\dist{a_n}{\ell} < 1$

According to the \xref[triangle inequality]{lec03::lem:triangle-inequality-distance} for distance, and $\dist{a_n}{\ell} < 1$, we have
\[\dist{0}{a_n} \leq \dist{0}{\ell} + \dist{\ell}{a_n} < \dist{0}{\ell} + 1\] 

Which means for all $n \in \NNOne$ with $n > N$, $\dist{0}{a_n} < \dist{0}{\ell} + 1$

And for $n \in \NNOne$ with $n \leq N$, we have $\dist{0}{a_n} \leq \max\{\dist{0}{a_1}, \dist{0}{a_2}, ..., \dist{0}{a_N}\}$

So, we can take $M = \max\{\max\{\dist{0}{a_1}, \dist{0}{a_2}, ..., \dist{0}{a_N}\}, \dist{0}{\ell} + 1\}$. Hence $M$ exists for all convergent sequences $a$
\end{proof} 

\begin{rem}
Although Convergent $\implies$ Bounded, the converse is not true. A bounded sequence may not be convergent

For example, the sequence $a_n = (-1)^n$ is bounded (since $\dist{a_n}{0} = 1$ for all $n$), but not convergent
\end{rem}

\section{Monotone Convergence Theorem} 

\begin{dfn}[Increasing Sequence]{increasing-sequence}
  A sequence $a: \NNOne \to F$ is said to be \emph{increasing} if
  \[\forall n \in \NNOne, \term{a}{n} \leq \term{a}{n+1}\]
\end{dfn}

\begin{dfn}[Strictly Increasing Sequence]{strictly-increasing-sequence}
  A sequence $a: \NNOne \to F$ is said to be \emph{strictly increasing} if
  \[\forall n \in \NNOne, \term{a}{n} < \term{a}{n+1}\]
\end{dfn} 

\begin{dfn}[Decreasing Sequence]{decreasing-sequence}
  A sequence $a: \NNOne \to F$ is said to be \emph{decreasing} if
  \[\forall n \in \NNOne, \term{a}{n} \geq \term{a}{n+1}\]
\end{dfn}

\begin{dfn}[Strictly Decreasing Sequence]{strictly-decreasing-sequence}
  A sequence $a: \NNOne \to F$ is said to be \emph{strictly decreasing} if
  \[\forall n \in \NNOne, \term{a}{n} > \term{a}{n+1}\]
\end{dfn}

\begin{lem}[Strictly Increasing Extends to All Indices]{strictly-increasing-extends}
  Let $a: \NNOne \to F$ be \xref[strictly increasing]{def:strictly-increasing-sequence}. Then $\forall n, m \in \NNOne,\, n > m \implies \term{a}{n} > \term{a}{m}$
\end{lem}

\begin{lem}[Increasing Extends to All Indices]{increasing-extends}
  Let $a: \NNOne \to F$ be \xref[increasing]{def:increasing-sequence}. Then $\forall n, m \in \NNOne,\, n \geq m \implies \term{a}{n} \geq \term{a}{m}$
\end{lem}

\begin{lem}[Strictly Decreasing Extends to All Indices]{strictly-decreasing-extends}
  Let $a: \NNOne \to F$ be \xref[strictly decreasing]{def:strictly-decreasing-sequence}. Then $\forall n, m \in \NNOne,\, n > m \implies \term{a}{n} < \term{a}{m}$
\end{lem}

\begin{lem}[Decreasing Extends to All Indices]{decreasing-extends}
  Let $a: \NNOne \to F$ be \xref[decreasing]{def:decreasing-sequence}. Then $\forall n, m \in \NNOne,\, n \geq m \implies \term{a}{n} \leq \term{a}{m}$
\end{lem}

\begin{dfn}[Monotone Sequence]{monotone-sequence}
  A sequence $a: \NNOne \to F$ is said to be \emph{monotone} if it is either \xref[increasing]{def:increasing-sequence} or \xref[decreasing]{def:decreasing-sequence}
\end{dfn} 

\begin{thm}[Monotone Convergence Theorem]{mct}
  Let \(F\) be a \xref[complete ordered field]{lec01::def:complete-ordered-field},
  and let \(a:\Bdd(F)\) be a \xref[monotone sequence]{def:monotone-sequence}.
  Then \(a\) is \xref[convergent]{def:convergent-sequence}
\end{thm}

\begin{proof}
Suppose $a$ is a \xref[bounded sequence]{def:bounded-sequence} in $F$ and \xref[increasing sequence]{def:increasing-sequence}, which means exist $x = \sup\{\term{a}{n}:n\in\NNOne\}$ (by the \xref[Axiom of Completeness]{lec01::axm:aoc}, applied to the nonempty set $\{\term{a}{n} : n \in \NNOne\}$ bounded above in $F$), and $\forall n, m \in \NNOne, \, n \leq m \implies \term{a}{n} \leq \term{a}{m} $

According to the \xref[$\varepsilon$ criterion for supremum]{lec01::prop:epsilon-supremum}, $\forall n \in \NNOne, \, \term{a}{n} \leq x $ and $\forall \varepsilon \in \apos{F}, \, \exists N \in \NNOne, \, \term{a}{N} > x - \varepsilon$

Adding the increasing property, ie \[N \leq n \implies \term{a}{N} \leq \term{a}{n} \leq x \implies (\term{a}{N} > x - \varepsilon \implies \term{a}{n} \geq \term{a}{N} > x - \varepsilon)\]

We can have \[\forall \varepsilon \in \apos{F}, \, \exists N \in \NNOne, \, \forall n \in \NNOne \geq N, \, \term{a}{n} > x - \varepsilon\]

Also, \[\term{a}{n} \leq x \, \land\, \varepsilon > 0 \implies \term{a}{n} < x + \varepsilon \]
\[\term{a}{n} > x - \varepsilon \, \land\,  \term{a}{n} < x + \varepsilon \implies x - \varepsilon < \term{a}{n} < x + \varepsilon \implies \dist{\term{a}{n}}{x} < \varepsilon\]
Notice that such $x$ is the \xref[supremum]{lec01::def:supremum} of the upper bounded sequence $a$, it exists. 
\[\exists x \in F, \, \forall \varepsilon \in \apos{F}, \, \exists N \in \NNOne, \, \forall n \in \NNOne \geq N, \, \dist{\term{a}{n}}{x} < \varepsilon\]

Which means it is a convergent sequence and its limit is $x$

If instead $a$ is \xref[decreasing]{def:decreasing-sequence}, the same argument applies with $x \defeq \inf\{\term{a}{n}:n\in\NNOne\}$ (which exists by the \xref[Axiom of Completeness]{lec01::axm:aoc}, since $\{\term{a}{n} : n \in \NNOne\}$ is nonempty and bounded below), using the dual $\varepsilon$-criterion for the infimum
\end{proof}

\section{The Number \texorpdfstring{\(e\)}{e}}

\begin{dfn}[Defining Sequence for \(e\)]{e-defining-sequence}
\end{dfn}

\begin{thm}[Convergence of the Defining Sequence for \(e\)]{e-sequence-convergence}
\end{thm}

\begin{dfn}[The Number \(e\)]{euler-number}
\end{dfn}

\section{Comparison of Sequences}

\begin{thm}[Comparison for Monotone Sequences]{monotone-comparison-increasing}
  Let $F$ be a \xref[complete ordered field]{lec01::def:complete-ordered-field}, and let $s,t : \NNOne \to F$ be \xref[increasing]{def:increasing-sequence} with $\term{s}{n} \leq \term{t}{n}$ for all $n \in \NNOne$. If $t$ is \xref[convergent]{def:convergent-sequence}, then
  \[
    s \text{ is convergent and }\seqlim(s) \leq \seqlim(t)
  \]
\end{thm}

\begin{proof}
\todo
\end{proof}

\begin{thm}[Comparison for Decreasing Sequences]{monotone-comparison-decreasing}
  Let $F$ be a \xref[complete ordered field]{lec01::def:complete-ordered-field}, and let $s,t : \NNOne \to F$ be \xref[decreasing]{def:decreasing-sequence} with $\term{s}{n} \leq \term{t}{n}$ for all $n \in \NNOne$. If $s$ is \xref[convergent]{def:convergent-sequence}, then
  \[
    t \text{ is convergent and }\seqlim(s) \leq \seqlim(t)
  \]
\end{thm}

\begin{proof}
Apply \xref{thm:monotone-comparison-increasing} to the increasing sequences $n\map-\term{t}{n}$ and $n\map-\term{s}{n}$
\end{proof}

\section{Squeeze Theorem}

\begin{thm}[Squeeze Theorem]{squeeze-theorem}
  Let $F$ be an \xref[ordered field]{lec01::def:ordered-field}, and let $a, x, b : \NNOne \to F$ with $\term{a}{n} \leq \term{x}{n} \leq \term{b}{n}$ for all $n \in \NNOne$. If $a$ and $b$ are \xref[convergent]{def:convergent-sequence} with $\seqlim(a) = \seqlim(b) = L$, then $x$ is convergent and $\seqlim(x) = L$
\end{thm}

\begin{proof}
\todo
\end{proof}

\end{document}
